% !TeX program = xelatex
\documentclass[12pt,a4paper]{article}
\usepackage[margin=2cm]{geometry}
\usepackage{fontspec}
\usepackage{amsmath,amssymb}
\usepackage{enumitem}
\setmainfont{Times New Roman}
\setlength{\parindent}{0pt}
\setlength{\parskip}{6pt}
\begin{document}
\begin{center}
{\Large\bfseries ĐỀ THI TUYỂN SINH VÀO LỚP 10 THPT}\\
{\bfseries Thành phố Hà Nội, năm học 2009--2010}\\
{\itshape Phần Đại số}
\end{center}

\textbf{Câu 1.} Cho biểu thức
\[
A=\frac{x}{x-4}+\frac{1}{\sqrt{x}-2}+\frac{1}{\sqrt{x}+2}
\qquad (x\ge0,\ x\ne4).
\]
\begin{enumerate}[label=\alph*)]
\item Rút gọn biểu thức $A$.
\item Tính giá trị của $A$ khi $x=25$.
\item Tìm giá trị của $x$ để $A=-\dfrac13$.
\end{enumerate}

\textbf{Câu 2.} Giải bài toán sau bằng cách lập phương trình hoặc hệ phương trình:

Hai tổ sản xuất cùng may một loại áo. Nếu tổ thứ nhất may trong 3 ngày, tổ thứ hai may trong 5 ngày thì cả hai tổ may được 1310 chiếc áo. Biết rằng trong một ngày tổ thứ nhất may được nhiều hơn tổ thứ hai 10 chiếc áo. Hỏi mỗi tổ trong một ngày may được bao nhiêu chiếc áo?

\textbf{Câu 3.} Cho phương trình
\[
x^2-2(m+1)x+m^2+2=0. \tag{1}
\]
\begin{enumerate}[label=\alph*)]
\item Giải phương trình khi $m=1$.
\item Tìm $m$ để phương trình có hai nghiệm phân biệt $x_1,x_2$ thỏa mãn
\[
x_1^2+x_2^2=10.
\]
\end{enumerate}

\textbf{Câu 5.} Giải phương trình
\[
x^2-\frac14+\sqrt{x^2+x+\frac14}
=\frac12\left(2x^3+x^2+2x+1\right).
\]
\end{document}
